\documentclass[aspectratio=169,10pt]{beamer}
\usetheme{Madrid}
\usecolortheme{default}
\usepackage{amsmath,amssymb}
\usepackage{enumitem}
\setlist[enumerate]{itemsep=1pt,topsep=1pt,leftmargin=*}

\title{Databricks ML Associate\\100 Q\&A Rapid Review}
\author{Exam Prep Deck}
\date{}

\begin{document}
	
	\frame{\titlepage}
	
	% ---------- TEMPLATE NOTE ----------
	\begin{frame}{How to use this deck}
		\begin{itemize}
			\item Odd slide: the exam-style question.
			\item Next slide: correct answer \textbf{+ why the other options are wrong}.
			\item Recommendation: cover the answer slide, answer aloud, then reveal.
		\end{itemize}
	\end{frame}
	
	% ============================================================
	% DOMAIN 1: DATABRICKS MACHINE LEARNING
	% ============================================================
	
	\begin{frame}{Q1 — MLflow Tracking}
		A data scientist needs parameters, metrics, and the model artifact logged automatically for a scikit-learn experiment. Best approach?
		\begin{enumerate}[label=\Alph*.]
			\item Manually write values to a CSV.
			\item Use \texttt{mlflow.autolog()} or explicit \texttt{log\_param}/\texttt{log\_metric} within a run.
			\item Print values to the notebook output.
			\item Save a screenshot of the experiment.
		\end{enumerate}
	\end{frame}
	
	\begin{frame}{A1 — Answer: B}
		\textbf{Correct: B.}
		\begin{itemize}
			\item \texttt{mlflow.autolog()} captures params, metrics, and model artifacts for supported libraries in one line.
			\item A and D are manual and not reproducible.
			\item C is not tracked in the registry.
		\end{itemize}
	\end{frame}
	
	\begin{frame}{Q2 — Model Registry}
		A model must be promoted from development to production with a clear, auditable lifecycle. Recommended Databricks approach?
		\begin{enumerate}[label=\Alph*.]
			\item Email the pickle file to the production team.
			\item Register in Unity Catalog with versions and an alias such as \texttt{@champion}.
			\item Rename the artifact folder.
			\item Copy the model file to a shared DBFS path.
		\end{enumerate}
	\end{frame}
	
	\begin{frame}{A2 — Answer: B}
		\textbf{Correct: B.}
		\begin{itemize}
			\item Unity Catalog registry gives governance, lineage, versioning, and aliases.
			\item \texttt{@champion} is the modern replacement for old "Production" stage.
			\item A, C, D have no audit trail and no governance.
		\end{itemize}
	\end{frame}
	
	\begin{frame}{Q3 — Feature Store}
		A team wants to avoid training-serving skew. Which Databricks capability is designed for this?
		\begin{enumerate}[label=\Alph*.]
			\item A shared Excel spreadsheet of feature definitions.
			\item Unity Catalog Feature Store with point-in-time joins.
			\item Recomputing features separately in serving code.
			\item Using only raw columns at inference.
		\end{enumerate}
	\end{frame}
	
	\begin{frame}{A3 — Answer: B}
		\textbf{Correct: B.}
		\begin{itemize}
			\item Feature Store = single source of truth for features used in training and serving.
			\item Point-in-time joins prevent future-data leakage.
			\item C is exactly what creates training-serving skew.
		\end{itemize}
	\end{frame}
	
	\begin{frame}{Q4 — AutoML}
		A data scientist wants a quick baseline on a labeled tabular dataset. Most appropriate feature?
		\begin{enumerate}[label=\Alph*.]
			\item Databricks AutoML.
			\item Building a custom neural network from scratch.
			\item Deploying an endpoint before training.
			\item Creating an online feature table only.
		\end{enumerate}
	\end{frame}
	
	\begin{frame}{A4 — Answer: A}
		\textbf{Correct: A.}
		\begin{itemize}
			\item AutoML produces baseline models, leaderboard, and generated notebooks.
			\item B is slow; C is out of order; D doesn't train anything.
		\end{itemize}
	\end{frame}
	
	\begin{frame}{Q5 — ML Runtime}
		Primary advantage of Databricks Runtime for Machine Learning?
		\begin{enumerate}[label=\Alph*.]
			\item Includes common ML libraries and integrations preinstalled.
			\item Eliminates all package conflicts automatically.
			\item Removes the need for cluster permissions.
			\item Allows skipping feature engineering entirely.
		\end{enumerate}
	\end{frame}
	
	\begin{frame}{A5 — Answer: A}
		\textbf{Correct: A.}
		\begin{itemize}
			\item ML Runtime ships with MLflow, Hyperopt, sklearn, XGBoost, etc.
			\item B is false — conflicts can still occur.
			\item C and D are obviously wrong.
		\end{itemize}
	\end{frame}
	
	\begin{frame}{Q6 — Unity Catalog Registry}
		Key benefit of registering models in Unity Catalog vs workspace registry?
		\begin{enumerate}[label=\Alph*.]
			\item Models become accessible only to the original author.
			\item Models are governed, discoverable, and promotable across workspaces at account level.
			\item Models no longer need a version number.
			\item Models are automatically deployed to serving endpoints.
		\end{enumerate}
	\end{frame}
	
	\begin{frame}{A6 — Answer: B}
		\textbf{Correct: B.}
		\begin{itemize}
			\item UC = account-level governance and cross-workspace discovery.
			\item A restricts access (opposite of intent).
			\item C is false — versioning always exists.
			\item D is false — deployment is separate.
		\end{itemize}
	\end{frame}
	
	\begin{frame}{Q7 — MLflow Search}
		How to programmatically find the best run by a metric?
		\begin{enumerate}[label=\Alph*.]
			\item Use \texttt{mlflow.search\_runs()} and sort the DataFrame.
			\item Manually click through each run in the UI.
			\item Export all runs to CSV and use Excel.
			\item Delete runs until one remains.
		\end{enumerate}
	\end{frame}
	
	\begin{frame}{A7 — Answer: A}
		\textbf{Correct: A.}
		\begin{itemize}
			\item \texttt{search\_runs()} returns a pandas DataFrame; sort by metric.
			\item B, C are manual; D is destructive.
		\end{itemize}
	\end{frame}
	
	\begin{frame}{Q8 — Model Aliases}
		What replaced older stages (Staging/Production/Archived) in Unity Catalog?
		\begin{enumerate}[label=\Alph*.]
			\item Numeric version numbers only.
			\item Aliases such as \texttt{champion} and \texttt{challenger}.
			\item Tags only.
			\item File system paths.
		\end{enumerate}
	\end{frame}
	
	\begin{frame}{A8 — Answer: B}
		\textbf{Correct: B.}
		\begin{itemize}
			\item Aliases are mutable pointers to a specific model version.
			\item E.g., \texttt{models:/my\_model@champion}.
			\item C (tags) are static metadata — not pointers.
		\end{itemize}
	\end{frame}
	
	\begin{frame}{Q9 — Online vs Offline Features}
		Primary difference between online and offline feature tables?
		\begin{enumerate}[label=\Alph*.]
			\item Online tables are for training; offline for serving.
			\item Online tables serve low-latency inference; offline for batch training/analysis.
			\item Online tables cannot be created in Unity Catalog.
			\item Offline tables do not support point-in-time joins.
		\end{enumerate}
	\end{frame}
	
	\begin{frame}{A9 — Answer: B}
		\textbf{Correct: B.}
		\begin{itemize}
			\item Online = low-latency lookup (e.g., DynamoDB-backed).
			\item Offline = large-scale batch / training reads.
			\item A reverses the roles; C and D are false.
		\end{itemize}
	\end{frame}
	
	\begin{frame}{Q10 — Feature Engineering Client}
		Which class creates/manages feature tables in Unity Catalog?
		\begin{enumerate}[label=\Alph*.]
			\item \texttt{FeatureEngineeringClient}.
			\item \texttt{SparkSession}.
			\item \texttt{MLflowClient}.
			\item \texttt{Pipeline}.
		\end{enumerate}
	\end{frame}
	
	\begin{frame}{A10 — Answer: A}
		\textbf{Correct: A.}
		\begin{itemize}
			\item \texttt{FeatureEngineeringClient} → \texttt{create\_table}, \texttt{write\_table}, \texttt{create\_training\_set}.
			\item \texttt{MLflowClient} manages runs/registry, not feature tables.
		\end{itemize}
	\end{frame}
	
	\begin{frame}{Q11 — AutoML Output}
		What useful artifact does AutoML generate for review and customization?
		\begin{enumerate}[label=\Alph*.]
			\item A screenshot of the leaderboard only.
			\item Generated notebooks for data exploration and best trial.
			\item A compiled binary with no source code.
			\item A cluster event log.
		\end{enumerate}
	\end{frame}
	
	\begin{frame}{A11 — Answer: B}
		\textbf{Correct: B.}
		\begin{itemize}
			\item AutoML generates a data exploration notebook and a best-trial training notebook.
			\item These can be opened, edited, and re-run as a starting point.
		\end{itemize}
	\end{frame}
	
	\begin{frame}{Q12 — MLOps Best Practice}
		Which setup best reflects sound MLOps strategy?
		\begin{enumerate}[label=\Alph*.]
			\item One shared workspace where notebooks are edited directly in production.
			\item Separate dev/staging/prod environments with versioned code and tracked artifacts.
			\item Manual model file copies between users.
			\item Training models without keeping prior run metadata.
		\end{enumerate}
	\end{frame}
	
	\begin{frame}{A12 — Answer: B}
		\textbf{Correct: B.}
		\begin{itemize}
			\item Environment separation + version control + experiment tracking = repeatable, auditable.
			\item A is the classic "everyone edits prod" antipattern.
		\end{itemize}
	\end{frame}
	
	\begin{frame}{Q13 — Reproducibility}
		Which practice most directly improves reproducibility?
		\begin{enumerate}[label=\Alph*.]
			\item Keeping training code in source control and logging params/artifacts in MLflow.
			\item Relying on notebook revision history alone.
			\item Letting each data scientist rename columns locally.
			\item Saving only final predictions.
		\end{enumerate}
	\end{frame}
	
	\begin{frame}{A13 — Answer: A}
		\textbf{Correct: A.}
		\begin{itemize}
			\item Reproducibility = code + params + artifacts all pinned and logged.
			\item Notebook history alone misses data \& environment state.
		\end{itemize}
	\end{frame}
	
	\begin{frame}{Q14 — Model Promotion Strategy}
		When is promoting code preferred over promoting models?
		\begin{enumerate}[label=\Alph*.]
			\item When the model is a simple scikit-learn model.
			\item When the training pipeline is complex and the model artifact is small.
			\item When the model artifact is large and training is expensive.
			\item When the team uses only AutoML.
		\end{enumerate}
	\end{frame}
	
	\begin{frame}{A14 — Answer: B}
		\textbf{Correct: B.}
		\begin{itemize}
			\item Promote-code: rebuild the model in each environment from versioned code — cleaner for complex pipelines with cheap training.
			\item Promote-model: ship the artifact — better when training is expensive (C).
		\end{itemize}
	\end{frame}
	
	\begin{frame}{Q15 — MLflow UI}
		Where in the MLflow UI can a user find the notebook that was run to create a run?
		\begin{enumerate}[label=\Alph*.]
			\item The ``Source'' link for the run.
			\item The ``Artifacts'' tab only.
			\item The cluster event log.
			\item The Databricks home page.
		\end{enumerate}
	\end{frame}
	
	\begin{frame}{A15 — Answer: A}
		\textbf{Correct: A.}
		\begin{itemize}
			\item ``Source'' links back to the notebook and revision used for that run.
			\item Artifacts tab only shows stored files.
		\end{itemize}
	\end{frame}
	
	\begin{frame}{Q16 — Feature Store Creation}
		Which method creates a feature table in Unity Catalog?
		\begin{enumerate}[label=\Alph*.]
			\item \texttt{fe.create\_table()}.
			\item \texttt{spark.sql("CREATE TABLE")}.
			\item \texttt{mlflow.log\_artifact()}.
			\item \texttt{dbutils.fs.mkdirs()}.
		\end{enumerate}
	\end{frame}
	
	\begin{frame}{A16 — Answer: A}
		\textbf{Correct: A.}
		\begin{itemize}
			\item \texttt{FeatureEngineeringClient.create\_table()} registers schema, primary keys, timestamp keys.
			\item Plain \texttt{CREATE TABLE} does not register as a feature table.
		\end{itemize}
	\end{frame}
	
	\begin{frame}{Q17 — Point-in-Time Lookup}
		Why is point-in-time lookup important for feature-store training sets?
		\begin{enumerate}[label=\Alph*.]
			\item Ensures training does not use feature values unavailable at prediction time.
			\item It speeds up the query.
			\item It automatically imputes missing values.
			\item It converts categorical features to numeric.
		\end{enumerate}
	\end{frame}
	
	\begin{frame}{A17 — Answer: A}
		\textbf{Correct: A.}
		\begin{itemize}
			\item Point-in-time join fetches each entity's feature values as of the label timestamp.
			\item Prevents label leakage from future data.
		\end{itemize}
	\end{frame}
	
	\begin{frame}{Q18 — Model Tags}
		How can a tag be set on a registered model in Unity Catalog?
		\begin{enumerate}[label=\Alph*.]
			\item \texttt{mlflow.set\_model\_tag()}.
			\item \texttt{spark.conf.set()}.
			\item \texttt{dbutils.notebook.exit()}.
			\item \texttt{display()}.
		\end{enumerate}
	\end{frame}
	
	\begin{frame}{A18 — Answer: A}
		\textbf{Correct: A.}
		\begin{itemize}
			\item \texttt{mlflow.set\_model\_tag()} and \texttt{mlflow.set\_model\_version\_tag()} handle model/version tags.
			\item Others are unrelated APIs.
		\end{itemize}
	\end{frame}
	
	\begin{frame}{Q19 — MLflow Run Identification}
		How to identify the best run using the MLflow Client API?
		\begin{enumerate}[label=\Alph*.]
			\item \texttt{client.search\_runs()}.
			\item \texttt{client.list\_experiments()}.
			\item \texttt{client.get\_run()}.
			\item \texttt{client.delete\_run()}.
		\end{enumerate}
	\end{frame}
	
	\begin{frame}{A19 — Answer: A}
		\textbf{Correct: A.}
		\begin{itemize}
			\item \texttt{search\_runs()} supports filters like \texttt{metrics.accuracy > 0.9} and ordering.
			\item \texttt{get\_run()} needs a run\_id you already know.
		\end{itemize}
	\end{frame}
	
	\begin{frame}{Q20 — Champion Promotion}
		How is a challenger promoted to champion using aliases?
		\begin{enumerate}[label=\Alph*.]
			\item Delete the old model and upload the new one.
			\item Move the \texttt{@champion} alias to the new model version.
			\item Rename the model.
			\item Create a new experiment.
		\end{enumerate}
	\end{frame}
	
	\begin{frame}{A20 — Answer: B}
		\textbf{Correct: B.}
		\begin{itemize}
			\item Aliases are mutable pointers — reassign \texttt{@champion} to new version.
			\item This is zero-downtime: consumers keep using \texttt{@champion}.
		\end{itemize}
	\end{frame}
	
	\begin{frame}{Q21 — MLflow Experiment Organization}
		A data scientist wants to track hyperparameter tuning with many child runs. Best MLflow pattern?
		\begin{enumerate}[label=\Alph*.]
			\item Separate experiment for each trial.
			\item Start child runs inside the parent run using \texttt{mlflow.start\_run(nested=True)}.
			\item Log all trials to one run without nesting.
			\item Avoid logging parameters.
		\end{enumerate}
	\end{frame}
	
	\begin{frame}{A21 — Answer: B}
		\textbf{Correct: B.}
		\begin{itemize}
			\item Nested runs group all hyperparameter trials under one parent.
			\item Easier to compare trials and identify the best config.
		\end{itemize}
	\end{frame}
	
	\begin{frame}{Q22 — Workspace vs UC Registry}
		Limitation of the workspace MLflow Model Registry vs Unity Catalog?
		\begin{enumerate}[label=\Alph*.]
			\item Workspace models cannot be versioned.
			\item Workspace models are scoped to a single workspace and lack account-level governance.
			\item Workspace models cannot be loaded in notebooks.
			\item Workspace models cannot be logged with MLflow.
		\end{enumerate}
	\end{frame}
	
	\begin{frame}{A22 — Answer: B}
		\textbf{Correct: B.}
		\begin{itemize}
			\item Workspace registry lives inside one workspace — no cross-workspace promotion.
			\item UC registry is account-wide and governed.
		\end{itemize}
	\end{frame}
	
	\begin{frame}{Q23 — MLflow Model Loading}
		Which function loads a logged MLflow model for inference?
		\begin{enumerate}[label=\Alph*.]
			\item \texttt{mlflow.pyfunc.load\_model()}.
			\item \texttt{mlflow.register\_model()}.
			\item \texttt{mlflow.log\_model()}.
			\item \texttt{mlflow.start\_run()}.
		\end{enumerate}
	\end{frame}
	
	\begin{frame}{A23 — Answer: A}
		\textbf{Correct: A.}
		\begin{itemize}
			\item \texttt{mlflow.pyfunc.load\_model()} or flavor-specific \texttt{load\_model()} loads for scoring.
			\item \texttt{log\_model()} logs, \texttt{register\_model()} registers — neither loads for scoring.
		\end{itemize}
	\end{frame}
	
	\begin{frame}{Q24 — Online Table Sync}
		How are online feature tables kept synchronized with offline tables?
		\begin{enumerate}[label=\Alph*.]
			\item Manual export/import.
			\item Databricks automatically publishes online tables from offline tables.
			\item Online tables are never synchronized.
			\item Only batch jobs can synchronize them.
		\end{enumerate}
	\end{frame}
	
	\begin{frame}{A24 — Answer: B}
		\textbf{Correct: B.}
		\begin{itemize}
			\item Databricks automatically syncs offline → online.
			\item This is part of the Feature Store design to reduce serving latency.
		\end{itemize}
	\end{frame}
	
	\begin{frame}{Q25 — AutoML Constraints}
		Which constraint can be specified when running Databricks AutoML?
		\begin{enumerate}[label=\Alph*.]
			\item A time limit for the experiment.
			\item A specific random seed only.
			\item A requirement to use only neural networks.
			\item A fixed number of features.
		\end{enumerate}
	\end{frame}
	
	\begin{frame}{A25 — Answer: A}
		\textbf{Correct: A.}
		\begin{itemize}
			\item \texttt{timeout\_minutes} caps total AutoML runtime.
			\item \texttt{exclude\_frameworks} can exclude specific algorithms — but no "only NN" option.
		\end{itemize}
	\end{frame}
	
	\begin{frame}{Q26 — MLflow Run Deletion}
		Which API call deletes an MLflow run?
		\begin{enumerate}[label=\Alph*.]
			\item \texttt{client.delete\_run(run\_id)}.
			\item \texttt{client.get\_run(run\_id)}.
			\item \texttt{client.search\_runs()}.
			\item \texttt{mlflow.end\_run()}.
		\end{enumerate}
	\end{frame}
	
	\begin{frame}{A26 — Answer: A}
		\textbf{Correct: A.}
		\begin{itemize}
			\item \texttt{delete\_run()} marks a run as deleted.
			\item \texttt{end\_run()} only terminates an active run — it does not delete it.
		\end{itemize}
	\end{frame}
	
	\begin{frame}{Q27 — Production Candidate Marking}
		Recommended way to mark a model version as current production candidate in UC?
		\begin{enumerate}[label=\Alph*.]
			\item Set the \texttt{@champion} alias on that version.
			\item Delete all other versions.
			\item Rename the model to include ``prod''.
			\item Create a new schema.
		\end{enumerate}
	\end{frame}
	
	\begin{frame}{A27 — Answer: A}
		\textbf{Correct: A.}
		\begin{itemize}
			\item \texttt{@champion} is the intended mechanism.
			\item Deleting versions destroys history; renaming is not portable.
		\end{itemize}
	\end{frame}
	
	\begin{frame}{Q28 — Experiment Artifacts}
		Where are MLflow run artifacts stored by default in Databricks?
		\begin{enumerate}[label=\Alph*.]
			\item In local cluster disk only.
			\item In the MLflow experiment's artifact location (DBFS or cloud storage).
			\item In notebook revision history.
			\item In the driver's memory.
		\end{enumerate}
	\end{frame}
	
	\begin{frame}{A28 — Answer: B}
		\textbf{Correct: B.}
		\begin{itemize}
			\item Each experiment has an \texttt{artifact\_location} (DBFS or cloud URI).
			\item Local disk is ephemeral; revision history stores notebook code.
		\end{itemize}
	\end{frame}
	
	\begin{frame}{Q29 — Registry Permissions}
		What permission is required to promote a model version to champion in UC?
		\begin{enumerate}[label=\Alph*.]
			\item \texttt{USE CATALOG} only.
			\item \texttt{EXECUTE} on the model.
			\item \texttt{MANAGE} or \texttt{APPLY TAG} on the model.
			\item \texttt{SELECT} on the training table.
		\end{enumerate}
	\end{frame}
	
	\begin{frame}{A29 — Answer: C}
		\textbf{Correct: C.}
		\begin{itemize}
			\item Aliases/tags are metadata edits → need \texttt{MANAGE} or \texttt{APPLY TAG}.
			\item \texttt{SELECT} on the training table doesn't grant model metadata rights.
		\end{itemize}
	\end{frame}
	
	\begin{frame}{Q30 — Autologging Scope}
		Which libraries does \texttt{mlflow.autolog()} support in Databricks?
		\begin{enumerate}[label=\Alph*.]
			\item Only scikit-learn.
			\item Multiple including scikit-learn, XGBoost, LightGBM, and Spark ML.
			\item Only TensorFlow.
			\item Only PyTorch.
		\end{enumerate}
	\end{frame}
	
	\begin{frame}{A30 — Answer: B}
		\textbf{Correct: B.}
		\begin{itemize}
			\item Databricks autologging supports sklearn, XGBoost, LightGBM, Spark ML, TensorFlow, PyTorch, and more.
		\end{itemize}
	\end{frame}
	
	\begin{frame}{Q31 — Feature Store Lookup}
		Which method retrieves features from a feature table for inference?
		\begin{enumerate}[label=\Alph*.]
			\item \texttt{fe.get\_online\_features()}.
			\item \texttt{spark.read.table()} only.
			\item \texttt{mlflow.load\_model()}.
			\item \texttt{dbutils.fs.head()}.
		\end{enumerate}
	\end{frame}
	
	\begin{frame}{A31 — Answer: A}
		\textbf{Correct: A.}
		\begin{itemize}
			\item \texttt{get\_online\_features()} returns low-latency feature records for a list of entity IDs.
			\item \texttt{spark.read.table()} reads batch features — not for real-time.
		\end{itemize}
	\end{frame}
	
	\begin{frame}{Q32 — Model Signature Enforcement}
		What happens when a model with signature is served and input does not match?
		\begin{enumerate}[label=\Alph*.]
			\item The request fails validation.
			\item The model silently ignores the mismatch.
			\item The model retrains automatically.
			\item The signature is updated.
		\end{enumerate}
	\end{frame}
	
	\begin{frame}{A32 — Answer: A}
		\textbf{Correct: A.}
		\begin{itemize}
			\item Signatures are enforced at inference; invalid schema → HTTP error.
			\item This is the whole point: catch serving-time schema drift.
		\end{itemize}
	\end{frame}
	
	\begin{frame}{Q33 — Tracking Server Location}
		In Databricks, where are MLflow experiments and runs stored?
		\begin{enumerate}[label=\Alph*.]
			\item On driver's local disk.
			\item In the Databricks-managed MLflow tracking service.
			\item In an external Git repository.
			\item In the cluster event log.
		\end{enumerate}
	\end{frame}
	
	\begin{frame}{A33 — Answer: B}
		\textbf{Correct: B.}
		\begin{itemize}
			\item Databricks provides a managed tracking server — no local setup needed.
			\item Workspace is the default tracking server.
		\end{itemize}
	\end{frame}
	
	\begin{frame}{Q34 — Deployment Readiness}
		Which check is most important before promoting to production?
		\begin{enumerate}[label=\Alph*.]
			\item Only training accuracy.
			\item Performance on a held-out test set and validation of input schema.
			\item Number of notebook cells.
			\item Cluster size used for training.
		\end{enumerate}
	\end{frame}
	
	\begin{frame}{A34 — Answer: B}
		\textbf{Correct: B.}
		\begin{itemize}
			\item Production readiness = good generalization + correct input schema.
			\item Training accuracy alone overestimates performance.
		\end{itemize}
	\end{frame}
	
	\begin{frame}{Q35 — Feature Table Primary Keys}
		Purpose of primary keys in a UC feature table?
		\begin{enumerate}[label=\Alph*.]
			\item Uniquely identify each entity for point-in-time lookups.
			\item Encrypt the table.
			\item Reduce storage costs.
			\item Replace the timestamp column.
		\end{enumerate}
	\end{frame}
	
	\begin{frame}{A35 — Answer: A}
		\textbf{Correct: A.}
		\begin{itemize}
			\item Primary keys + timestamp keys = join keys for PIT lookups.
			\item They do not encrypt or replace timestamps.
		\end{itemize}
	\end{frame}
	
	\begin{frame}{Q36 — MLflow Run Hierarchy}
		Which structure correctly represents experiments/runs/artifacts?
		\begin{enumerate}[label=\Alph*.]
			\item Experiments contain runs; runs contain artifacts.
			\item Runs contain experiments; experiments contain artifacts.
			\item Artifacts contain runs; runs contain experiments.
			\item Experiments and runs are unrelated.
		\end{enumerate}
	\end{frame}
	
	\begin{frame}{A36 — Answer: A}
		\textbf{Correct: A.}
		\begin{itemize}
			\item Experiment → many runs → each run logs artifacts, params, metrics.
		\end{itemize}
	\end{frame}
	
	\begin{frame}{Q37 — AutoML Explainability}
		What does AutoML provide to help interpret the best model?
		\begin{enumerate}[label=\Alph*.]
			\item SHAP-based feature importance and generated notebooks.
			\item Only a leaderboard.
			\item Only final predictions.
			\item A compiled binary.
		\end{enumerate}
	\end{frame}
	
	\begin{frame}{A37 — Answer: A}
		\textbf{Correct: A.}
		\begin{itemize}
			\item AutoML leaderboard provides SHAP summaries and per-row explanations.
			\item Notebooks make the winning pipeline reproducible.
		\end{itemize}
	\end{frame}
	
	\begin{frame}{Q38 — Registry Lineage}
		What information does Unity Catalog retain for lineage?
		\begin{enumerate}[label=\Alph*.]
			\item Only the model name.
			\item Run ID, source experiment, and model version lineage.
			\item Only file size.
			\item Only upload timestamp.
		\end{enumerate}
	\end{frame}
	
	\begin{frame}{A38 — Answer: B}
		\textbf{Correct: B.}
		\begin{itemize}
			\item Lineage links model version → run → experiment → input tables.
			\item This enables audit and impact analysis.
		\end{itemize}
	\end{frame}
	
	% ============================================================
	% DOMAIN 2: MODEL DEVELOPMENT
	% ============================================================
	
	\begin{frame}{Q39 — Regression Metrics}
		RMSE = 25, MAE = 20, $R^2$ = 0.85. Which two statements are correct?
		\begin{enumerate}[label=\Alph*.]
			\item Model captures 85\% of variance.
			\item RMSE > MAE suggests larger errors exist.
			\item Model is underfitting.
			\item MAE is 25.
		\end{enumerate}
	\end{frame}
	
	\begin{frame}{A39 — Answer: A and B}
		\textbf{Correct: A, B.}
		\begin{itemize}
			\item \$R\^2 = 0.85\$ → 85\% of variance explained.
			\item RMSE penalizes large errors more than MAE  RMSE > MAE  some large errors.
			\item D is wrong — MAE is 20, not 25.
		\end{itemize}
	\end{frame}
	
	\begin{frame}{Q40 — Imbalanced Classification}
		98\% accuracy, 2\% positives. What should be done?
		\begin{enumerate}[label=\Alph*.]
			\item Ship immediately.
			\item Evaluate precision, recall, F1, PR-AUC; consider class weights or resampling (training only).
			\item Increase epochs.
			\item Remove positive class.
		\end{enumerate}
	\end{frame}
	
	\begin{frame}{A40 — Answer: B}
		\textbf{Correct: B.}
		\begin{itemize}
			\item Accuracy is misleading on imbalanced classes.
			\item Use PR-AUC, F1, recall; resample or reweight on train only.
			\item Never resample test set (would inflate metrics).
		\end{itemize}
	\end{frame}
	
	\begin{frame}{Q41 — One-Hot Encoding}
		When is one-hot encoding least appropriate?
		\begin{enumerate}[label=\Alph*.]
			\item Low cardinality.
			\item Extremely high cardinality.
			\item Using linear models.
			\item Feature is a string.
		\end{enumerate}
	\end{frame}
	
	\begin{frame}{A41 — Answer: B}
		\textbf{Correct: B.}
		\begin{itemize}
			\item High cardinality → too many sparse columns → memory blow-up.
			\item Prefer target encoding or feature hashing.
		\end{itemize}
	\end{frame}
	
	\begin{frame}{Q42 — Imputation}
		For a highly skewed distribution, preferred imputation strategy?
		\begin{enumerate}[label=\Alph*.]
			\item Mean.
			\item Median.
			\item Mode.
			\item Zero.
		\end{enumerate}
	\end{frame}
	
	\begin{frame}{A42 — Answer: B}
		\textbf{Correct: B.}
		\begin{itemize}
			\item Median is robust to skew and outliers.
			\item Mean gets pulled by skew and large values.
		\end{itemize}
	\end{frame}
	
	\begin{frame}{Q43 — Log Transformation}
		When is a log transformation most appropriate?
		\begin{enumerate}[label=\Alph*.]
			\item Normally distributed feature.
			\item Right-skewed feature spanning orders of magnitude.
			\item Categorical feature.
			\item Binary target.
		\end{enumerate}
	\end{frame}
	
	\begin{frame}{A43 — Answer: B}
		\textbf{Correct: B.}
		\begin{itemize}
			\item Log compresses large values, reduces skew, stabilizes variance.
			\item Not meaningful for already-normal or categorical features.
		\end{itemize}
	\end{frame}
	
	\begin{frame}{Q44 — Estimator vs Transformer}
		Key difference between an Estimator and a Transformer in Spark ML?
		\begin{enumerate}[label=\Alph*.]
			\item Estimator has \texttt{fit()}; Transformer has \texttt{transform()}.
			\item Estimator transforms data; Transformer fits models.
			\item They are identical.
			\item Transformer cannot be used in a Pipeline.
		\end{enumerate}
	\end{frame}
	
	\begin{frame}{A44 — Answer: A}
		\textbf{Correct: A.}
		\begin{itemize}
			\item Estimator.fit(df) → Transformer.
			\item Transformer.transform(df) → df.
			\item Both are Pipeline stages.
		\end{itemize}
	\end{frame}
	
	\begin{frame}{Q45 — Spark ML Pipeline}
		Why wrap preprocessing + training in a Spark ML Pipeline?
		\begin{enumerate}[label=\Alph*.]
			\item Ensures same preprocessing at inference as training.
			\item Reduces the number of columns.
			\item Avoids using MLflow.
			\item Eliminates the need for a test set.
		\end{enumerate}
	\end{frame}
	
	\begin{frame}{A45 — Answer: A}
		\textbf{Correct: A.}
		\begin{itemize}
			\item PipelineModel captures fitted preprocessing parameters.
			\item Prevents train/serve skew.
		\end{itemize}
	\end{frame}
	
	\begin{frame}{Q46 — Hyperopt Parallelization}
		Which Hyperopt component distributes trials across a Spark cluster?
		\begin{enumerate}[label=\Alph*.]
			\item \texttt{SparkTrials}.
			\item \texttt{Trials}.
			\item \texttt{fmin}.
			\item \texttt{hp.uniform}.
		\end{enumerate}
	\end{frame}
	
	\begin{frame}{A46 — Answer: A}
		\textbf{Correct: A.}
		\begin{itemize}
			\item \texttt{SparkTrials} → runs trials in parallel on Spark workers.
			\item \texttt{Trials} is single-node serial.
		\end{itemize}
	\end{frame}
	
	\begin{frame}{Q47 — Hyperopt Algorithm}
		Which algorithm does Hyperopt use via \texttt{tpe.suggest}?
		\begin{enumerate}[label=\Alph*.]
			\item Grid search.
			\item Random search.
			\item Tree-structured Parzen Estimator (Bayesian optimization).
			\item Exhaustive search.
		\end{enumerate}
	\end{frame}
	
	\begin{frame}{A47 — Answer: C}
		\textbf{Correct: C.}
		\begin{itemize}
			\item TPE = sequential model-based (Bayesian) optimizer.
			\item Smarter than grid/random at the same budget.
		\end{itemize}
	\end{frame}
	
	\begin{frame}{Q48 — CV vs Single Split}
		Main trade-off between cross-validation and a single train-validation split?
		\begin{enumerate}[label=\Alph*.]
			\item CV is faster.
			\item CV gives a more stable estimate but is more expensive.
			\item Single split always more accurate.
			\item Same cost.
		\end{enumerate}
	\end{frame}
	
	\begin{frame}{A48 — Answer: B}
		\textbf{Correct: B.}
		\begin{itemize}
			\item $k$-fold CV = $k$ model fits → more robust but slower.
			\item Single split is noisier but cheap.
		\end{itemize}
	\end{frame}
	
	\begin{frame}{Q49 — Imbalanced Metrics}
		Best metric for a highly imbalanced binary classification task?
		\begin{enumerate}[label=\Alph*.]
			\item Accuracy.
			\item PR-AUC.
			\item MAE.
			\item $R^2$.
		\end{enumerate}
	\end{frame}
	
	\begin{frame}{A49 — Answer: B}
		\textbf{Correct: B.}
		\begin{itemize}
			\item PR-AUC focuses on minority class performance.
			\item ROC-AUC can look optimistic under heavy imbalance.
		\end{itemize}
	\end{frame}
	
	\begin{frame}{Q50 — Random Split}
		Which Spark operation randomly splits a DataFrame into train and test?
		\begin{enumerate}[label=\Alph*.]
			\item \texttt{DataFrame.randomSplit()}.
			\item \texttt{TrainValidationSplit}.
			\item \texttt{CrossValidator}.
			\item \texttt{DataFrame.where()}.
		\end{enumerate}
	\end{frame}
	
	\begin{frame}{A50 — Answer: A}
		\textbf{Correct: A.}
		\begin{itemize}
			\item \texttt{df.randomSplit([0.8, 0.2], seed=42)} → list of DataFrames.
			\item \texttt{TrainValidationSplit} is for hyperparameter tuning, not raw splitting.
		\end{itemize}
	\end{frame}
	
	\begin{frame}{Q51 — StringIndexer}
		Before one-hot encoding a string column in Spark ML, which stage is typically required?
		\begin{enumerate}[label=\Alph*.]
			\item \texttt{StringIndexer}.
			\item \texttt{VectorAssembler}.
			\item \texttt{StandardScaler}.
			\item \texttt{Imputer}.
		\end{enumerate}
	\end{frame}
	
	\begin{frame}{A51 — Answer: A}
		\textbf{Correct: A.}
		\begin{itemize}
			\item \texttt{OneHotEncoder} needs numeric index input, produced by \texttt{StringIndexer}.
			\item Typical chain: StringIndexer → OneHotEncoder → VectorAssembler.
		\end{itemize}
	\end{frame}
	
	\begin{frame}{Q52 — VectorAssembler}
		Purpose of \texttt{VectorAssembler}?
		\begin{enumerate}[label=\Alph*.]
			\item Combine multiple feature columns into a vector column.
			\item Impute missing values.
			\item Encode categorical variables.
			\item Split data.
		\end{enumerate}
	\end{frame}
	
	\begin{frame}{A52 — Answer: A}
		\textbf{Correct: A.}
		\begin{itemize}
			\item Spark ML models require a single "features" vector column.
			\item \texttt{VectorAssembler} does that final assembly step.
		\end{itemize}
	\end{frame}
	
	\begin{frame}{Q53 — Feature Scaling}
		Which scaler makes features zero-mean/unit-variance?
		\begin{enumerate}[label=\Alph*.]
			\item \texttt{MinMaxScaler}.
			\item \texttt{StandardScaler}.
			\item \texttt{MaxAbsScaler}.
			\item \texttt{Normalizer}.
		\end{enumerate}
	\end{frame}
	
	\begin{frame}{A53 — Answer: B}
		\textbf{Correct: B.}
		\begin{itemize}
			\item \texttt{StandardScaler}: $(x - \mu)/\sigma$.
			\item \texttt{MinMaxScaler} → $[0,1]$; \texttt{Normalizer} → unit L-p norm row-wise.
		\end{itemize}
	\end{frame}
	
	\begin{frame}{Q54 — Data Leakage}
		Standardizing before train/test split. What is the risk?
		\begin{enumerate}[label=\Alph*.]
			\item Test statistics leak into train → overly optimistic evaluation.
			\item Underfitting.
			\item Features become categorical.
			\item Nothing; best practice.
		\end{enumerate}
	\end{frame}
	
	\begin{frame}{A54 — Answer: A}
		\textbf{Correct: A.}
		\begin{itemize}
			\item Scaler should be fit on train only, then applied to test.
			\item Otherwise leakage inflates metrics.
		\end{itemize}
	\end{frame}
	
	\begin{frame}{Q55 — Pandas UDF}
		Benefit of vectorized pandas UDFs over standard PySpark UDFs?
		\begin{enumerate}[label=\Alph*.]
			\item Batch processing (Arrow) → much faster.
			\item Only work on single-node DataFrames.
			\item Require no Python.
			\item Auto-handle missing values.
		\end{enumerate}
	\end{frame}
	
	\begin{frame}{A55 — Answer: A}
		\textbf{Correct: A.}
		\begin{itemize}
			\item Vectorized pandas UDFs use Apache Arrow for batch transfer → big speedup.
			\item Standard UDFs are row-by-row (slow).
		\end{itemize}
	\end{frame}
	
	\begin{frame}{Q56 — Target Encoding}
		Target encoding is most useful for which situation?
		\begin{enumerate}[label=\Alph*.]
			\item Low-cardinality categorical.
			\item High-cardinality categorical.
			\item Continuous with outliers.
			\item Binary targets only.
		\end{enumerate}
	\end{frame}
	
	\begin{frame}{A56 — Answer: B}
		\textbf{Correct: B.}
		\begin{itemize}
			\item Replaces category with mean target — compact for many categories.
			\item Use CV/leave-one-out to avoid leakage.
		\end{itemize}
	\end{frame}
	
	\begin{frame}{Q57 — Overfitting}
		Good train performance, poor test performance. What is this called?
		\begin{enumerate}[label=\Alph*.]
			\item Underfitting.
			\item Overfitting.
			\item Data leakage.
			\item Regularization.
		\end{enumerate}
	\end{frame}
	
	\begin{frame}{A57 — Answer: B}
		\textbf{Correct: B.}
		\begin{itemize}
			\item Model memorized training noise → poor generalization.
			\item Fix: regularization, more data, simpler model.
		\end{itemize}
	\end{frame}
	
	\begin{frame}{Q58 — Bias-Variance}
		High variance typically corresponds to which problem?
		\begin{enumerate}[label=\Alph*.]
			\item Overfitting.
			\item Underfitting.
			\item Data imbalance.
			\item Feature scaling.
		\end{enumerate}
	\end{frame}
	
	\begin{frame}{A58 — Answer: A}
		\textbf{Correct: A.}
		\begin{itemize}
			\item High variance = very sensitive to training sample = overfit.
			\item High bias = too simple = underfit.
		\end{itemize}
	\end{frame}
	
	\begin{frame}{Q59 — Median vs Mean Imputation}
		When is median imputation preferred?
		\begin{enumerate}[label=\Alph*.]
			\item Symmetric feature.
			\item Feature with outliers/skew.
			\item Categorical feature.
			\item No missing values.
		\end{enumerate}
	\end{frame}
	
	\begin{frame}{A59 — Answer: B}
		\textbf{Correct: B.}
		\begin{itemize}
			\item Median is robust to skew/outliers.
			\item For symmetric, mean is fine.
		\end{itemize}
	\end{frame}
	
	\begin{frame}{Q60 — Class Weights}
		How to mitigate imbalance without resampling?
		\begin{enumerate}[label=\Alph*.]
			\item Class weights in training.
			\item Remove minority samples.
			\item Duplicate majority class.
			\item Ignore the problem.
		\end{enumerate}
	\end{frame}
	
	\begin{frame}{A60 — Answer: A}
		\textbf{Correct: A.}
		\begin{itemize}
			\item \texttt{weightCol} or \texttt{class\_weight} upweights minority class.
			\item B, C make the problem worse.
		\end{itemize}
	\end{frame}
	
	\begin{frame}{Q61 — Metric Selection}
		Large errors are especially costly in a regression task. Best metric?
		\begin{enumerate}[label=\Alph*.]
			\item RMSE.
			\item MAE.
			\item Accuracy.
			\item F1.
		\end{enumerate}
	\end{frame}
	
	\begin{frame}{A61 — Answer: A}
		\textbf{Correct: A.}
		\begin{itemize}
			\item RMSE squares errors → penalizes large deviations.
			\item MAE treats all errors linearly.
		\end{itemize}
	\end{frame}
	
	\begin{frame}{Q62 — Spark ML Persistence}
		How is a fitted Spark ML model saved?
		\begin{enumerate}[label=\Alph*.]
			\item \texttt{model.write().overwrite().save(path)}.
			\item \texttt{mlflow.log\_artifact()}.
			\item \texttt{spark.sql("SAVE")}.
			\item \texttt{dbutils.fs.put()}.
		\end{enumerate}
	\end{frame}
	
	\begin{frame}{A62 — Answer: A}
		\textbf{Correct: A.}
		\begin{itemize}
			\item Spark MLModel has \texttt{save()}/\texttt{write()} — saves directory with metadata + Parquet.
			\item \texttt{log\_artifact()} logs to MLflow — different mechanism.
		\end{itemize}
	\end{frame}
	
	\begin{frame}{Q63 — CrossValidator Grid}
		How is the hyperparameter grid specified in Spark ML \texttt{CrossValidator}?
		\begin{enumerate}[label=\Alph*.]
			\item \texttt{ParamGridBuilder}.
			\item Python dict only.
			\item CSV.
			\item SQL query.
		\end{enumerate}
	\end{frame}
	
	\begin{frame}{A63 — Answer: A}
		\textbf{Correct: A.}
		\begin{itemize}
			\item \texttt{ParamGridBuilder().addGrid(...).build()}.
		\end{itemize}
	\end{frame}
	
	\begin{frame}{Q64 — Evaluator Choice}
		Which evaluator pairs with \texttt{CrossValidator} for regression?
		\begin{enumerate}[label=\Alph*.]
			\item \texttt{RegressionEvaluator}.
			\item \texttt{BinaryClassificationEvaluator}.
			\item \texttt{MulticlassClassificationEvaluator}.
			\item \texttt{ClusteringEvaluator}.
		\end{enumerate}
	\end{frame}
	
	\begin{frame}{A64 — Answer: A}
		\textbf{Correct: A.}
		\begin{itemize}
			\item \texttt{RegressionEvaluator} supports RMSE, MAE, $R^2$.
			\item Classification evaluators expect label/score structure.
		\end{itemize}
	\end{frame}
	
	\begin{frame}{Q65 — Hyperopt Objective}
		What must the Hyperopt objective return?
		\begin{enumerate}[label=\Alph*.]
			\item Scalar loss to minimize.
			\item DataFrame.
			\item Trained model only.
			\item String.
		\end{enumerate}
	\end{frame}
	
	\begin{frame}{A65 — Answer: A}
		\textbf{Correct: A.}
		\begin{itemize}
			\item \texttt{fmin(fn, space, algo, max\_evals)} → returns lowest scalar loss.
			\item To maximize, return \texttt{-metric}.
		\end{itemize}
	\end{frame}
	
	\begin{frame}{Q66 — Hyperopt Space}
		Uniform distribution over a range?
		\begin{enumerate}[label=\Alph*.]
			\item \texttt{hp.uniform('x', low, high)}.
			\item \texttt{hp.choice('x', [a,b])}.
			\item \texttt{hp.quniform('x', low, high, q)}.
			\item \texttt{hp.loguniform('x', low, high)}.
		\end{enumerate}
	\end{frame}
	
	\begin{frame}{A66 — Answer: A}
		\textbf{Correct: A.}
		\begin{itemize}
			\item \texttt{hp.uniform} = continuous uniform.
			\item \texttt{quniform} = quantized, \texttt{loguniform} = log-scale.
		\end{itemize}
	\end{frame}
	
	\begin{frame}{Q67 — Logistic Regression Class}
		Which Spark ML class trains binary logistic regression?
		\begin{enumerate}[label=\Alph*.]
			\item \texttt{LogisticRegression}.
			\item \texttt{LinearRegression}.
			\item \texttt{DecisionTreeClassifier}.
			\item \texttt{RandomForestRegressor}.
		\end{enumerate}
	\end{frame}
	
	\begin{frame}{A67 — Answer: A}
		\textbf{Correct: A.}
	\end{frame}
	
	\begin{frame}{Q68 — Linear Regression Class}
		Which Spark ML class trains linear regression?
		\begin{enumerate}[label=\Alph*.]
			\item \texttt{LinearRegression}.
			\item \texttt{LogisticRegression}.
			\item \texttt{DecisionTreeClassifier}.
			\item \texttt{NaiveBayes}.
		\end{enumerate}
	\end{frame}
	
	\begin{frame}{A68 — Answer: A}
		\textbf{Correct: A.}
	\end{frame}
	
	\begin{frame}{Q69 — Feature Importance}
		How is feature importance accessed for a tree model?
		\begin{enumerate}[label=\Alph*.]
			\item \texttt{model.featureImportances}.
			\item \texttt{model.coefficients}.
			\item \texttt{model.intercept}.
			\item \texttt{model.summary}.
		\end{enumerate}
	\end{frame}
	
	\begin{frame}{A69 — Answer: A}
		\textbf{Correct: A.}
		\begin{itemize}
			\item Tree-based models expose \texttt{featureImportances} (vector).
			\item Linear models use \texttt{coefficients}.
		\end{itemize}
	\end{frame}
	
	\begin{frame}{Q70 — Linear Coefficients}
		How are coefficients accessed for Spark ML linear regression?
		\begin{enumerate}[label=\Alph*.]
			\item \texttt{model.coefficients}.
			\item \texttt{model.featureImportances}.
			\item \texttt{model.intercept} only.
			\item \texttt{model.summary} only.
		\end{enumerate}
	\end{frame}
	
	\begin{frame}{A70 — Answer: A}
		\textbf{Correct: A.}
	\end{frame}
	
	\begin{frame}{Q71 — StringIndexer Invalid}
		Default behavior for unseen label at transform time?
		\begin{enumerate}[label=\Alph*.]
			\item Throws error unless \texttt{handleInvalid} set.
			\item Silently ignores row.
			\item Maps to index 0.
			\item Maps to most frequent.
		\end{enumerate}
	\end{frame}
	
	\begin{frame}{A71 — Answer: A}
		\textbf{Correct: A.}
		\begin{itemize}
			\item Default = \texttt{"error"}.
			\item \texttt{handleInvalid="skip"} drops row; \texttt{"keep"} → index = numLabels.
		\end{itemize}
	\end{frame}
	
	\begin{frame}{Q72 — OneHotEncoder Input}
		What input does \texttt{OneHotEncoder} expect?
		\begin{enumerate}[label=\Alph*.]
			\item Numeric index column (from StringIndexer).
			\item Raw string column.
			\item Continuous feature.
			\item Vector column.
		\end{enumerate}
	\end{frame}
	
	\begin{frame}{A72 — Answer: A}
		\textbf{Correct: A.}
		\begin{itemize}
			\item \texttt{StringIndexer} → \texttt{OneHotEncoder} → \texttt{VectorAssembler}.
		\end{itemize}
	\end{frame}
	
	\begin{frame}{Q73 — VectorAssembler Input Types}
		What column types can \texttt{VectorAssembler} accept?
		\begin{enumerate}[label=\Alph*.]
			\item Numeric, boolean, and vector columns.
			\item Only string columns.
			\item Only categorical columns.
			\item Only binary columns.
		\end{enumerate}
	\end{frame}
	
	\begin{frame}{A73 — Answer: A}
	\end{frame}
	
	\begin{frame}{Q74 — Pipeline Fit Output}
		What does fitting a \texttt{Pipeline} produce?
		\begin{enumerate}[label=\Alph*.]
			\item \texttt{PipelineModel}.
			\item Raw DataFrame.
			\item Cluster.
			\item MLflow experiment.
		\end{enumerate}
	\end{frame}
	
	\begin{frame}{A74 — Answer: A}
	\end{frame}
	
	\begin{frame}{Q75 — Model Save Format}
		What format does Spark ML use to save models?
		\begin{enumerate}[label=\Alph*.]
			\item Directory with metadata and Parquet.
			\item Single CSV.
			\item Pickle only.
			\item JSON only.
		\end{enumerate}
	\end{frame}
	
	\begin{frame}{A75 — Answer: A}
	\end{frame}
	
	\begin{frame}{Q76 — MLflow Model Flavor}
		What is an MLflow model flavor?
		\begin{enumerate}[label=\Alph*.]
			\item Convention for saving/loading models from different libraries.
			\item Type of cluster.
			\item Visualization tool.
			\item Data source.
		\end{enumerate}
	\end{frame}
	
	\begin{frame}{A76 — Answer: A}
	\end{frame}
	
	\begin{frame}{Q77 — Model Signature}
		Purpose of a model signature?
		\begin{enumerate}[label=\Alph*.]
			\item Defines expected input/output schema.
			\item Encrypts the model.
			\item Compresses the model.
			\item Deletes the model.
		\end{enumerate}
	\end{frame}
	
	\begin{frame}{A77 — Answer: A}
	\end{frame}
	
	\begin{frame}{Q78 — StandardScaler Params}
		Which parameters does \texttt{StandardScaler} use?
		\begin{enumerate}[label=\Alph*.]
			\item Mean and std.
			\item Min and max.
			\item Median and IQR.
			\item Mode and variance.
		\end{enumerate}
	\end{frame}
	
	\begin{frame}{A78 — Answer: A}
	\end{frame}
	
	\begin{frame}{Q79 — MinMaxScaler Range}
		Default scaling range of \texttt{MinMaxScaler}?
		\begin{enumerate}[label=\Alph*.]
			\item $[0,1]$.
			\item $[-1,1]$.
			\item $[0,100]$.
			\item Mean 0, var 1.
		\end{enumerate}
	\end{frame}
	
	\begin{frame}{A79 — Answer: A}
	\end{frame}
	
	\begin{frame}{Q80 — Imputer Strategies}
		Which strategies does Spark ML \texttt{Imputer} support?
		\begin{enumerate}[label=\Alph*.]
			\item Mean, median, mode.
			\item Only mean.
			\item Only zero.
			\item Only median.
		\end{enumerate}
	\end{frame}
	
	\begin{frame}{A80 — Answer: A}
	\end{frame}
	
	\begin{frame}{Q81 — CV numFolds}
		What does \texttt{numFolds} control in \texttt{CrossValidator}?
		\begin{enumerate}[label=\Alph*.]
			\item Number of $k$ folds.
			\item Number of models to deploy.
			\item Number of features to select.
			\item Number of clusters.
		\end{enumerate}
	\end{frame}
	
	\begin{frame}{A81 — Answer: A}
	\end{frame}
	
	\begin{frame}{Q82 — trainRatio}
		What does \texttt{trainRatio} control in \texttt{TrainValidationSplit}?
		\begin{enumerate}[label=\Alph*.]
			\item Fraction used for training.
			\item Number of folds.
			\item Number of trials.
			\item Learning rate.
		\end{enumerate}
	\end{frame}
	
	\begin{frame}{A82 — Answer: A}
	\end{frame}
	
	\begin{frame}{Q83 — DecisionTreeClassifier}
		Which class trains a decision tree classifier?
		\begin{enumerate}[label=\Alph*.]
			\item \texttt{DecisionTreeClassifier}.
			\item \texttt{LinearRegression}.
			\item \texttt{LogisticRegression}.
			\item \texttt{KMeans}.
		\end{enumerate}
	\end{frame}
	
	\begin{frame}{A83 — Answer: A}
	\end{frame}
	
	\begin{frame}{Q84 — RandomForestClassifier}
		Which class trains a random forest classifier?
		\begin{enumerate}[label=\Alph*.]
			\item \texttt{RandomForestClassifier}.
			\item \texttt{DecisionTreeRegressor}.
			\item \texttt{LinearSVC}.
			\item \texttt{NaiveBayes}.
		\end{enumerate}
	\end{frame}
	
	\begin{frame}{A84 — Answer: A}
	\end{frame}
	
	\begin{frame}{Q85 — GBT}
		Which class trains a gradient-boosted tree classifier?
		\begin{enumerate}[label=\Alph*.]
			\item \texttt{GBTClassifier}.
			\item \texttt{RandomForestClassifier}.
			\item \texttt{DecisionTreeClassifier}.
			\item \texttt{LogisticRegression}.
		\end{enumerate}
	\end{frame}
	
	\begin{frame}{A85 — Answer: A}
	\end{frame}
	
	\begin{frame}{Q86 — KMeans}
		Which class performs $k$-means clustering?
		\begin{enumerate}[label=\Alph*.]
			\item \texttt{KMeans}.
			\item \texttt{LogisticRegression}.
			\item \texttt{LinearRegression}.
			\item \texttt{RandomForestClassifier}.
		\end{enumerate}
	\end{frame}
	
	\begin{frame}{A86 — Answer: A}
	\end{frame}
	
	\begin{frame}{Q87 — Binary Metric}
		Which metric is available in \texttt{BinaryClassificationEvaluator}?
		\begin{enumerate}[label=\Alph*.]
			\item Area under ROC.
			\item RMSE.
			\item MAE.
			\item $R^2$.
		\end{enumerate}
	\end{frame}
	
	\begin{frame}{A87 — Answer: A}
	\end{frame}
	
	\begin{frame}{Q88 — Regression Metric}
		Which metric is available in \texttt{RegressionEvaluator}?
		\begin{enumerate}[label=\Alph*.]
			\item RMSE.
			\item Precision.
			\item Recall.
			\item F1.
		\end{enumerate}
	\end{frame}
	
	\begin{frame}{A88 — Answer: A}
	\end{frame}
	
	\begin{frame}{Q89 — Multiclass Metric}
		Which metric is available in \texttt{MulticlassClassificationEvaluator}?
		\begin{enumerate}[label=\Alph*.]
			\item F1.
			\item RMSE.
			\item MAE.
			\item $R^2$.
		\end{enumerate}
	\end{frame}
	
	\begin{frame}{A89 — Answer: A}
	\end{frame}
	
	\begin{frame}{Q90 — FeatureHasher}
		Purpose of \texttt{FeatureHasher}?
		\begin{enumerate}[label=\Alph*.]
			\item Hash categorical features to fixed-length vector without dictionary.
			\item Impute missing values.
			\item Scale features.
			\item Split data.
		\end{enumerate}
	\end{frame}
	
	\begin{frame}{A90 — Answer: A}
		\begin{itemize}
			\item No vocabulary needed → constant memory.
			\item Downside: hash collisions possible.
		\end{itemize}
	\end{frame}
	
	\begin{frame}{Q91 — SQLTransformer}
		Purpose of \texttt{SQLTransformer}?
		\begin{enumerate}[label=\Alph*.]
			\item Apply SQL expression to transform within a Pipeline.
			\item Train model.
			\item Evaluate model.
			\item Deploy model.
		\end{enumerate}
	\end{frame}
	
	\begin{frame}{A91 — Answer: A}
	\end{frame}
	
	\begin{frame}{Q92 — VectorSlicer}
		Purpose of \texttt{VectorSlicer}?
		\begin{enumerate}[label=\Alph*.]
			\item Select subset of features from a vector column.
			\item Combine columns into vector.
			\item Scale features.
			\item Impute missing.
		\end{enumerate}
	\end{frame}
	
	\begin{frame}{A92 — Answer: A}
	\end{frame}
	
	\begin{frame}{Q93 — PolynomialExpansion}
		Purpose of \texttt{PolynomialExpansion}?
		\begin{enumerate}[label=\Alph*.]
			\item Generate polynomial and interaction features.
			\item Encode categorical.
			\item Scale features.
			\item Split data.
		\end{enumerate}
	\end{frame}
	
	\begin{frame}{A93 — Answer: A}
	\end{frame}
	
	\begin{frame}{Q94 — Binarizer}
		Purpose of \texttt{Binarizer}?
		\begin{enumerate}[label=\Alph*.]
			\item Threshold a continuous feature into binary.
			\item Encode categorical.
			\item Scale features.
			\item Impute missing.
		\end{enumerate}
	\end{frame}
	
	\begin{frame}{A94 — Answer: A}
	\end{frame}
	
	% ============================================================
	% DOMAIN 3: ML WORKFLOWS
	% ============================================================
	
	\begin{frame}{Q95 — Summary Statistics}
		Which command computes summary statistics on a Spark DataFrame?
		\begin{enumerate}[label=\Alph*.]
			\item \texttt{df.summary()}.
			\item \texttt{df.show()}.
			\item \texttt{df.printSchema()}.
			\item \texttt{df.count()}.
		\end{enumerate}
	\end{frame}
	
	\begin{frame}{A95 — Answer: A}
		\begin{itemize}
			\item \texttt{df.summary()} → count/mean/std/min/quartiles/max for numeric columns.
			\item \texttt{describe()} gives a simpler set.
		\end{itemize}
	\end{frame}
	
	\begin{frame}{Q96 — Std-dev Outlier Removal}
		How to remove outliers based on standard deviation?
		\begin{enumerate}[label=\Alph*.]
			\item Filter values beyond $k\sigma$ from mean.
			\item Remove all above median.
			\item Remove top 50\%.
			\item \texttt{dropDuplicates()}.
		\end{enumerate}
	\end{frame}
	
	\begin{frame}{A96 — Answer: A}
		\begin{itemize}
			\item Common choice: $|x - \mu| > 3\sigma$ → drop.
			\item Assumes roughly normal distribution.
		\end{itemize}
	\end{frame}
	
	\begin{frame}{Q97 — IQR Method}
		IQR outlier detection rule?
		\begin{enumerate}[label=\Alph*.]
			\item Below Q1 or above Q3.
			\item $< Q1 - 1.5\,IQR$ or $> Q3 + 1.5\,IQR$.
			\item Beyond mean.
			\item $z > 1$.
		\end{enumerate}
	\end{frame}
	
	\begin{frame}{A97 — Answer: B}
		\begin{itemize}
			\item Standard Tukey fence: $Q1 - 1.5\,IQR$ to $Q3 + 1.5\,IQR$.
			\item More robust than $z$-score to non-normal data.
		\end{itemize}
	\end{frame}
	
	\begin{frame}{Q98 — Pandas API Refactor}
		Pandas notebook needs to scale to big data with minimal refactor. Best approach?
		\begin{enumerate}[label=\Alph*.]
			\item Rewrite in Scala.
			\item Refactor to pandas API on Spark.
			\item Use only Spark SQL.
			\item Sample to memory.
		\end{enumerate}
	\end{frame}
	
	\begin{frame}{A98 — Answer: B}
		\begin{itemize}
			\item pandas API on Spark mirrors pandas syntax on distributed data.
			\item Sampling loses information; Scala rewrite is expensive.
		\end{itemize}
	\end{frame}
	
	\begin{frame}{Q99 — Two Continuous Features}
		Which visualization for comparing two continuous features?
		\begin{enumerate}[label=\Alph*.]
			\item Scatter plot.
			\item Bar chart of counts.
			\item Pie chart.
			\item Box plot of one feature.
		\end{enumerate}
	\end{frame}
	
	\begin{frame}{A99 — Answer: A}
		\begin{itemize}
			\item Scatter shows relationship/correlation between two continuous variables.
		\end{itemize}
	\end{frame}
	
	\begin{frame}{Q100 — Two Categorical Features}
		Which visualization for comparing two categorical features?
		\begin{enumerate}[label=\Alph*.]
			\item Grouped bar chart or heatmap.
			\item Scatter plot.
			\item Line chart.
			\item Histogram.
		\end{enumerate}
	\end{frame}
	
	\begin{frame}{A100 — Answer: A}
		\begin{itemize}
			\item Crosstab → grouped bar or heatmap makes comparison easy.
		\end{itemize}
	\end{frame}
	
	\begin{frame}{Exam Day Tips}
		\begin{itemize}
			\item \textbf{Time budget}: ~1.9 min per question (48 Q / 90 min).
			\item Watch for \textbf{"Select two"} — these appear frequently.
			\item Kill time on traps: "accuracy" on imbalanced data, "scale before split", "means of serving" etc.
			\item \textbf{Priority domains}: Databricks ML (38%) + Model Development (31%) — master these first.
			\item Hands-on practice: MLflow + UC registry + Feature Store + Spark ML Pipeline.
		\end{itemize}
	\end{frame}
	
\end{document}